\documentclass[10pt,a4paper]{article}

% ---------- lingua, codifica, font (testo + matematica) ----------
\usepackage[T1]{fontenc}
\usepackage[utf8]{inputenc}
\usepackage[italian]{babel}
\usepackage{mathpazo}                 % Palatino per testo e matematica
\usepackage{microtype}

% ---------- matematica ----------
\usepackage{amsmath,amssymb,amsfonts,mathtools,bm,cancel,mathrsfs,nicefrac}
\DeclareMathOperator{\sen}{sen}       % seno all'italiana
\DeclareMathOperator{\tg}{tg}
\DeclareMathOperator{\ctg}{cotg}
\DeclareMathOperator{\arcsen}{arcsen}
\DeclareMathOperator{\arctg}{arctg}
\newcommand{\dd}{\mathrm{d}}                       % differenziale
\newcommand{\vv}[1]{\vec{#1}}                      % vettore con freccia
\newcommand{\norma}[1]{\left\lvert#1\right\rvert} % modulo di vettore
\newcommand{\pedice}[1]{_{\text{#1}}}              % pedice testuale: E\pedice{in}

% ---------- unità di misura e numeri (virgola decimale) ----------
\usepackage{siunitx}
\sisetup{output-decimal-marker={,},per-mode=symbol,
         inter-unit-product=\ensuremath{{}\cdot{}},group-digits=false}

% ---------- grafici e figure ----------
\usepackage{tikz}
\usetikzlibrary{arrows.meta,calc,positioning,patterns,decorations.pathmorphing}
\usepackage{pgfplots}
\pgfplotsset{compat=1.16}
\usepackage{graphicx}
% Grafici collocati nel punto pertinente dello svolgimento.
\newenvironment{grafico}{\begin{figure}[h!]\centering}{\end{figure}}

% ---------- impaginazione ----------
\usepackage[a4paper,left=1.8cm,right=1.8cm,top=2.3cm,bottom=2.2cm]{geometry}
\usepackage{fancyhdr,lastpage,needspace,enumitem,booktabs,array,colortbl,tabularx}
\setlength{\parindent}{0pt}
\setlength{\parskip}{3pt}
\setlength{\headheight}{14pt}
\setlist[enumerate]{leftmargin=1.7em,itemsep=2pt,topsep=2pt,parsep=0pt}
\setlist[itemize]{leftmargin=1.3em,itemsep=2pt,topsep=2pt,parsep=0pt}

% ---------- colori degli indicatori ----------
\usepackage[table,dvipsnames]{xcolor}
\definecolor{anac}{HTML}{C62828}\definecolor{anabg}{HTML}{FDECEA} % rosso
\definecolor{svic}{HTML}{1A56DB}\definecolor{svibg}{HTML}{E8F0FE} % blu
\definecolor{argc}{HTML}{A67C00}\definecolor{argbg}{HTML}{FFF6CC} % giallo
\definecolor{comc}{HTML}{2E7D32}\definecolor{combg}{HTML}{E6F4EA} % verde
% Solo griglie d'esame (su richiesta): IND arancio, DAT teal
\definecolor{indc}{HTML}{E65100}\definecolor{indbg}{HTML}{FFF0E0}
\definecolor{datc}{HTML}{00796B}\definecolor{datbg}{HTML}{E0F2F1}

% Riferimenti agli indicatori: usare queste macro in tutto il documento.
\newcommand{\coloreANA}{anac}\newcommand{\nomeANA}{Analizzare}
\newcommand{\coloreSVI}{svic}\newcommand{\nomeSVI}{Sviluppare e calcolare}
\newcommand{\coloreARG}{argc}\newcommand{\nomeARG}{Argomentare}
\newcommand{\coloreCOM}{comc}\newcommand{\nomeCOM}{Comunicare}
\newcommand{\coloreIND}{indc}\newcommand{\nomeIND}{Individuare}
\newcommand{\coloreDAT}{datc}\newcommand{\nomeDAT}{Interpretare i dati}
% Griglie d'esame: matematica \renewcommand{\nomeANA}{Comprendere}, \nomeSVI{Sviluppare il processo risolutivo};
% fisica/mista: \nomeSVI{Sviluppare il processo risolutivo}, \nomeDAT{Interpretare criticamente i dati} (mista: ...rappresentare, elaborare).
\newcommand{\sigla}[1]{\textcolor{\csname colore#1\endcsname}{\textbf{#1}}}
\newcommand{\indicatore}[1]{\textcolor{\csname colore#1\endcsname}{\textbf{#1: \csname nome#1\endcsname}}}

% ---------- box colorati e chip ----------
\usepackage{tcolorbox}
\tcbuselibrary{skins,breakable}

% Box passaggio chiave: \begin{boxana}[Titolo] ... \end{boxana}  (ana|svi|arg|com)
\newcommand{\defbox}[4]{%
  \newtcolorbox{box#1}[1][Passaggio chiave]{enhanced,breakable,frame hidden,
    borderline west={2.5pt}{0pt}{#2},colback=#3,boxrule=0pt,arc=1.5pt,
    left=7pt,right=5pt,top=2pt,bottom=2pt,before skip=4pt,after skip=4pt,
    fontupper=\small,before upper={\textbf{\idevidenza}\enspace\sigla{#4}\enspace\textbf{##1}\par}}}
\defbox{ana}{anac}{anabg}{ANA}
\defbox{svi}{svic}{svibg}{SVI}
\defbox{arg}{argc}{argbg}{ARG}
\defbox{com}{comc}{combg}{COM}
\defbox{ind}{indc}{indbg}{IND}
\defbox{dat}{datc}{datbg}{DAT}

% Box enunciato
\newtcolorbox{enun}{colback=black!3,colframe=black!25,boxrule=0.4pt,arc=2pt,
  left=6pt,right=6pt,top=3pt,bottom=3pt,fontupper=\small,
  before skip=5pt,after skip=5pt}

% Chip identificativo e massimo in punti; i dati sono forniti dal registro approvato.
\newcommand{\chip}[4]{\tcbox[on line,boxsep=0pt,left=3pt,right=3pt,top=1pt,bottom=1pt,
  arc=4pt,boxrule=0.4pt,colframe=#1,colback=#2,fontupper=\scriptsize\bfseries\color{#1}]{#3\ #4}}
\newcommand{\ANA}[2]{\chip{anac}{anabg}{\sigla{ANA}}{#1\, / massimo #2\,p}}
\newcommand{\SVI}[2]{\chip{svic}{svibg}{\sigla{SVI}}{#1\, / massimo #2\,p}}
\newcommand{\ARG}[2]{\chip{argc}{argbg}{\sigla{ARG}}{#1\, / massimo #2\,p}}
\newcommand{\COM}[2]{\chip{comc}{combg}{\sigla{COM}}{#1\, / massimo #2\,p}}
\newcommand{\IND}[2]{\chip{indc}{indbg}{\sigla{IND}}{#1\, / massimo #2\,p}}
\newcommand{\DAT}[2]{\chip{datc}{datbg}{\sigla{DAT}}{#1\, / massimo #2\,p}}
\newcommand{\ess}{\ensuremath{\bigstar}}
\newcommand{\legenda}{\indicatore{ANA}\quad\indicatore{SVI}\par
  \indicatore{ARG}\quad\indicatore{COM}}
% Griglie d'esame: \legendaesame{ANA}{IND}{SVI}{ARG} (matematica) o \legendaesame{ANA}{SVI}{DAT}{ARG} (fisica, mista)
\newcommand{\legendaesame}[4]{\indicatore{#1}\quad\indicatore{#2}\par\indicatore{#3}\quad\indicatore{#4}}
% Limite categorico sull'intera verifica; il registro non incrementa il conteggio.
\newcounter{numeroevidenze}
\newcommand{\controllaevidenze}{\stepcounter{numeroevidenze}%
  \ifnum\value{numeroevidenze}>40\relax
    \PackageError{correttore}{Superato il limite categorico di 40 evidenze}{Accorpare le evidenze prima di generare il correttore.}%
  \fi}
% Metadati mantenuti nel sorgente; nel corpo il titolo viene emesso solo dal box.
\newcommand{\idevidenza}{}
\newcommand{\evidenza}[5]{\controllaevidenze\def\idevidenza{#1}\needspace{8\baselineskip}}
% Registro raggruppato: esercizio, argomento comune, righe ID/nome/punto/indicatore/massimo.
\newcommand{\registrogruppo}[3]{\needspace{6\baselineskip}\par\smallskip{\small
  \textbf{Esercizio #1}\enspace #2\par
  \begin{tabularx}{\linewidth}{@{}lXllr@{}}\toprule
  ID & Evidenza & Punto & Ind. & Max\\\midrule
  #3\bottomrule\end{tabularx}\par}}
\newcommand{\inizioesercizi}{\clearpage}
\newcommand{\riepiloghi}[2]{\sez{Riepiloghi}
  \noindent\begin{minipage}[t]{0.58\linewidth}#1\end{minipage}\hfill
  \begin{minipage}[t]{0.39\linewidth}#2\end{minipage}\par}

\newcommand{\livelli}[5]{\begin{description}[leftmargin=2.7cm,style=nextline,font=\small\bfseries]
\item[Non raggiunto] #1
\if\relax\detokenize{#2}\relax\else\item[Iniziale] #2\fi
\item[Base] #3
\if\relax\detokenize{#4}\relax\else\item[Intermedio] #4\fi
\item[Avanzato] #5
\end{description}}

% Intestazioni di colonna colorate
\newcommand{\hANA}{\cellcolor{anabg}\sigla{ANA}}
\newcommand{\hSVI}{\cellcolor{svibg}\sigla{SVI}}
\newcommand{\hARG}{\cellcolor{argbg}\sigla{ARG}}
\newcommand{\hCOM}{\cellcolor{combg}\sigla{COM}}
\newcommand{\hIND}{\cellcolor{indbg}\sigla{IND}}
\newcommand{\hDAT}{\cellcolor{datbg}\sigla{DAT}}

% ---------- macro di struttura ----------
% Sezione semplice
\newcommand{\sez}[1]{\needspace{4\baselineskip}\par\vspace{1em}%
  {\large\bfseries #1}\par\vspace{-0.4em}\rule{\linewidth}{0.6pt}\par\vspace{0.2em}}
% Esercizio: \esercizio{Titolo}{ANA}{SVI}{ARG}{COM}{TOT}
\newcommand{\esercizio}[6]{\needspace{8\baselineskip}\par\vspace{1em}%
  {\large\bfseries #1}\par\vspace{0.25em}}
% Punto: \punto{a) titolo}
\newcommand{\punto}[1]{\needspace{6\baselineskip}\par\vspace{0.7em}{\bfseries #1}\par\vspace{0.2em}}
% Tabella punteggi del punto: \tabpunti{ANA}{SVI}{ARG}{COM}{TOT}  (usare 0 se assente)
\newcommand{\tabpunti}[5]{\par\noindent\renewcommand{\arraystretch}{1.15}\small
  \begin{tabular}{*{5}{|>{\centering\arraybackslash}p{1.4cm}}|}\hline
  \hANA & \hSVI & \hARG & \hCOM & \textbf{Tot}\\\hline
  #1 & #2 & #3 & #4 & \textbf{#5}\\\hline\end{tabular}\par\smallskip}
% Griglie d'esame: \tabpuntiesame{ANA}{IND}{SVI}{ARG}{v1}{v2}{v3}{v4}{TOT} (sigle nell'ordine dei valori)
\newcommand{\tabpuntiesame}[9]{\par\noindent\renewcommand{\arraystretch}{1.15}\small
  \begin{tabular}{*{5}{|>{\centering\arraybackslash}p{1.4cm}}|}\hline
  \csname h#1\endcsname & \csname h#2\endcsname & \csname h#3\endcsname & \csname h#4\endcsname & \textbf{Tot}\\\hline
  #5 & #6 & #7 & #8 & \textbf{#9}\\\hline\end{tabular}\par\smallskip}
% Nota (alternative, errori trascinati, criteri)
\newcommand{\nota}[1]{\begin{itemize}[font=\footnotesize]\item {\footnotesize #1}\end{itemize}}
\newcommand{\criteri}[1]{\par\nopagebreak[3]\smallskip\needspace{3\baselineskip}%
  {\footnotesize\textbf{Alternative, errori e criteri}\par
  \begin{itemize}[itemsep=1pt,topsep=1pt]#1\end{itemize}}}

% ---------- intestazione e piè di pagina ----------
\newcommand{\materia}{}
\newcommand{\classe}{}
\newcommand{\argomenti}{}
\pagestyle{fancy}\fancyhf{}
\renewcommand{\headrulewidth}{0.3pt}
\lhead{\small\textsc{Correttore\ \materia}}
\rhead{}
\cfoot{\small Pag.\ \thepage\ di \pageref{LastPage}}

\renewcommand{\materia}{Fisica}
\renewcommand{\classe}{Classe terza, liceo scientifico \textperiodcentered\ Urti e quantità di moto}
\begin{document}
\begin{center}
{\LARGE\bfseries Correttore\ \materia}\\[3pt]
\classe
\end{center}

\legenda
\sez{Registro delle evidenze}
\registrogruppo{1}{URT}{
E01 & Riconosce la conservazione della quantità di moto & intero & \sigla{ANA} & 40\\
E02 & Ricava e calcola la velocità finale & intero & \sigla{SVI} & 100\\
E03 & Calcola l'energia cinetica dissipata & intero & \sigla{SVI} & 80\\
E04 & Scrive formule, risultati e unità & intero & \sigla{COM} & 40\\}
\registrogruppo{2}{URT}{
E05 & Sceglie quantità di moto ed energia cinetica & intero & \sigla{ANA} & 50\\
E06 & Ricava le velocità finali & intero & \sigla{SVI} & 120\\
E07 & Commenta il segno e verifica l'energia cinetica & intero & \sigla{ARG} & 60\\}
\registrogruppo{3}{URT, ENE (indicati per evidenza)}{
E08 & Separa urto e risalita, ENE & intero & \sigla{ANA} & 70\\
E09 & Velocità del sistema dopo l'urto, ENE & intero & \sigla{SVI} & 110\\
E10 & Velocità del proiettile, URT & intero & \sigla{SVI} & 110\\
E11 & Spiega la non conservazione dell'energia, ENE & intero & \sigla{ARG} & 70\\
E12 & Organizza il percorso con unità, ENE & intero & \sigla{COM} & 80\\}
\registrogruppo{Intero elaborato}{TRASVERSALE, Competenze trasversali}{
E13 & Notazione, unità e cifre significative & Intero elaborato & \sigla{COM} & 70\\}

\riepiloghi{%
\begin{tabular}{@{}lrrrrr@{}}\toprule
 & \hANA & \hSVI & \hARG & \hCOM & Tot\\\midrule
Esercizio 1 & 40 & 180 & 0 & 40 & 260\\
Esercizio 2 & 50 & 120 & 60 & 0 & 230\\
Esercizio 3 & 70 & 220 & 70 & 80 & 440\\
Intero elaborato & 0 & 0 & 0 & 70 & 70\\\midrule
Totale & 160 & 520 & 130 & 190 & 1000\\\bottomrule
\end{tabular}}{%
\begin{tabular}{@{}lr@{}}\toprule
Argomento & Max\\\midrule
URT, Urti e quantità di moto & 600\\
ENE, Energia meccanica & 330\\
Trasversale & 70\\\midrule
Totale & 1000\\\bottomrule
\end{tabular}}

\inizioesercizi
\esercizio{Esercizio 1 -- Urto completamente anelastico}{}{}{}{}{}
\begin{enun}
Carrello $m_1=\SI{0.50}{kg}$ a $v_1=\SI{4.0}{m/s}$ contro $m_2=\SI{1.5}{kg}$ fermo; restano agganciati. Velocità finale ed energia dissipata.
\end{enun}
\evidenza{E01}{1}{ANA}{URT}{40}
\begin{boxana}[Riconosce la conservazione della quantità di moto]
\begin{enumerate}
\item Non agiscono forze esterne orizzontali sul sistema dei due carrelli: la quantità di moto totale si conserva.
\item L'energia cinetica non si conserva, perché l'urto è completamente anelastico.
\end{enumerate}
\livelli{Non individua la legge di conservazione adatta.}{}{Sceglie la conservazione della quantità di moto con la condizione sul sistema isolato.}{}{Aggiunge che l'energia cinetica non si conserva.}
\end{boxana}
\evidenza{E02}{1}{SVI}{URT}{100}
\begin{boxsvi}[Ricava e calcola la velocità finale]
\begin{enumerate}
\item $m_1v_1=(m_1+m_2)\,V$, da cui $V=\dfrac{m_1v_1}{m_1+m_2}$.
\item $V=\dfrac{\SI{0.50}{kg}\cdot\SI{4.0}{m/s}}{\SI{2.0}{kg}}=\SI{1.0}{m/s}$.
\end{enumerate}
\livelli{Non ricava la velocità finale.}{}{Ricava $V$ in lettere e sostituisce i dati, con eventuali imprecisioni.}{}{Ricava $V$ in lettere e la calcola con unità e verso corretti.}
\criteri{\item Ammessa la sostituzione diretta della formula del corpo unico purché derivata dalla conservazione.}
\end{boxsvi}
\evidenza{E03}{1}{SVI}{URT}{80}
\begin{boxsvi}[Calcola l'energia cinetica dissipata]
\begin{enumerate}
\item $K_i=\tfrac12 m_1v_1^2=\SI{4.0}{J}$.
\item $K_f=\tfrac12(m_1+m_2)V^2=\SI{1.0}{J}$.
\item $|\Delta K|=K_i-K_f=\SI{3.0}{J}$, il $75\%$ dell'energia iniziale.
\end{enumerate}
\livelli{Non calcola l'energia dissipata.}{}{Calcola $K_i$ e $K_f$ ma sbaglia la differenza o l'unità.}{}{Calcola $|\Delta K|=\SI{3.0}{J}$ con unità.}
\end{boxsvi}
\evidenza{E04}{1}{COM}{URT}{40}
\begin{boxcom}[Scrive formule, risultati e unità]
\begin{enumerate}
\item Le formule sono scritte in lettere prima della sostituzione dei dati.
\item Ogni risultato ha l'unità di misura al termine dell'espressione.
\end{enumerate}
\livelli{Formule e unità mancano.}{}{Usa formule e unità con qualche omissione.}{}{Formule in lettere e unità sempre corrette.}
\end{boxcom}

\esercizio{Esercizio 2 -- Urto elastico}{}{}{}{}{}
\begin{enun}
Blocco $m_1=\SI{2.0}{kg}$ a $v_1=\SI{4.0}{m/s}$ contro $m_2=\SI{3.0}{kg}$ fermo, urto elastico: velocità finali, segno e verifica dell'energia cinetica.
\end{enun}
\evidenza{E05}{2}{ANA}{URT}{50}
\begin{boxana}[Sceglie quantità di moto ed energia cinetica]
\begin{enumerate}
\item Nell'urto elastico si conservano la quantità di moto e l'energia cinetica totale.
\item Il sistema è: $m_1v_1=m_1u_1+m_2u_2$ e $\tfrac12m_1v_1^2=\tfrac12m_1u_1^2+\tfrac12m_2u_2^2$.
\end{enumerate}
\livelli{Sceglie una sola conservazione o nessuna.}{}{Scrive entrambe le leggi, con qualche imprecisione.}{}{Scrive correttamente il sistema delle due leggi.}
\end{boxana}
\evidenza{E06}{2}{SVI}{URT}{120}
\begin{boxsvi}[Ricava le velocità finali]
\begin{enumerate}
\item Dal sistema: $u_1=\dfrac{m_1-m_2}{m_1+m_2}\,v_1$ e $u_2=\dfrac{2m_1}{m_1+m_2}\,v_1$.
\item $u_1=\dfrac{\SI{2.0}{kg}-\SI{3.0}{kg}}{\SI{5.0}{kg}}\cdot\SI{4.0}{m/s}=\SI{-0.80}{m/s}$.
\item $u_2=\dfrac{2\cdot\SI{2.0}{kg}}{\SI{5.0}{kg}}\cdot\SI{4.0}{m/s}=\SI{3.2}{m/s}$.
\end{enumerate}
\livelli{Non ricava le velocità.}{}{Ricava una velocità o commette errori di segno.}{}{Ricava entrambe le velocità in lettere e con unità.}
\criteri{\item Ammessa la risoluzione del sistema per sostituzione o con la velocità relativa $u_2-u_1=v_1$.}
\end{boxsvi}
\evidenza{E07}{2}{ARG}{URT}{60}
\begin{boxarg}[Commenta il segno e verifica l'energia cinetica]
\begin{enumerate}
\item $u_1<0$: il blocco più leggero rimbalza e torna indietro.
\item $K_i=\SI{16}{J}$ e $K_f=\SI{0.64}{J}+\SI{15.36}{J}=\SI{16}{J}$: l'energia cinetica si conserva.
\end{enumerate}
\livelli{Nessun commento o verifica.}{}{Commenta il segno oppure verifica $K$, non entrambi.}{}{Commenta il segno e verifica la conservazione di $K$ con i valori.}
\end{boxarg}

\esercizio{Esercizio 3 -- Pendolo balistico}{}{}{}{}{}
\begin{enun}
Proiettile $m=\SI{20}{g}$ in un blocco $M=\SI{1.98}{kg}$; il sistema sale di $\SI{10.0}{cm}$. Velocità del proiettile e perché l'energia meccanica non si conserva nell'urto. Costante concessa: accelerazione di gravità $\SI{9.81}{m/s^2}$ (non valutata).
\end{enun}
\evidenza{E08}{3}{ANA}{ENE}{70}
\begin{boxana}[Separa urto e risalita]
\begin{enumerate}
\item L'urto è brevissimo e completamente anelastico: si conserva la quantità di moto.
\item Dopo l'urto agisce solo la gravità: si conserva l'energia meccanica del sistema.
\end{enumerate}
\livelli{Applica una sola legge a tutto il fenomeno.}{}{Distingue le due fasi con qualche confusione.}{}{Distingue le due fasi e assegna a ciascuna la legge corretta.}
\end{boxana}
\evidenza{E09}{3}{SVI}{ENE}{110}
\begin{boxsvi}[Ricava la velocità subito dopo l'urto]
\begin{enumerate}
\item $\tfrac12(m+M)V^2=(m+M)\,g\,h$, da cui $V=\sqrt{2gh}$.
\item $V=\sqrt{2\cdot\SI{9.81}{m/s^2}\cdot\SI{0.100}{m}}=\SI{1.40}{m/s}$.
\end{enumerate}
\livelli{Non determina $V$.}{}{Imposta l'energia ma sbaglia la massa o il calcolo.}{}{Ricava $V$ in lettere e la calcola con unità.}
\end{boxsvi}
\evidenza{E10}{3}{SVI}{URT}{110}
\begin{boxsvi}[Ricava la velocità del proiettile]
\begin{enumerate}
\item $mv=(m+M)\,V$, da cui $v=\dfrac{m+M}{m}\,V$.
\item $v=\dfrac{\SI{2.00}{kg}}{\SI{0.020}{kg}}\cdot\SI{1.40}{m/s}\approx\SI{1.4e2}{m/s}$.
\end{enumerate}
\livelli{Non ricava $v$.}{}{Ricava $v$ con errori di unità, di massa o di arrotondamento.}{}{Ricava $v$ in lettere e la calcola con unità e cifre significative corrette.}
\criteri{\item Un errore su $V$ è già penalizzato in E09: qui il seguito coerente vale il massimo.}
\end{boxsvi}
\evidenza{E11}{3}{ARG}{ENE}{70}
\begin{boxarg}[Spiega la non conservazione dell'energia]
\begin{enumerate}
\item Durante l'urto agiscono forze interne dissipative: parte dell'energia cinetica diventa energia interna.
\item Per questo l'energia meccanica si conserva solo dopo l'urto, non prima.
\end{enumerate}
\livelli{Nessuna spiegazione o spiegazione errata.}{}{Dice che l'urto è anelastico senza spiegare dove va l'energia.}{}{Spiega la dissipazione e distingue prima e dopo l'urto.}
\end{boxarg}
\evidenza{E12}{3}{COM}{ENE}{80}
\begin{boxcom}[Organizza il percorso con unità]
\begin{enumerate}
\item Le due fasi sono separate e le formule sono scritte in lettere prima di sostituire.
\item Le unità compaiono nei risultati intermedi e finali.
\end{enumerate}
\livelli{Percorso non leggibile.}{}{Percorso leggibile con unità mancanti.}{}{Percorso ordinato in due fasi, con formule in lettere e unità.}
\end{boxcom}

\esercizio{Intero elaborato}{}{}{}{}{}
\evidenza{E13}{Intero elaborato}{COM}{TRASVERSALE}{70}
\begin{boxcom}[Usa notazione, unità e cifre significative]
\begin{enumerate}
\item I simboli distinguono velocità prima e dopo l'urto, con pedici chiari.
\item Le cifre significative dei risultati sono coerenti con i dati.
\end{enumerate}
\livelli{Notazione confusa.}{}{Notazione corretta con cifre significative non sempre scelte.}{}{Notazione e cifre significative sempre corrette.}
\criteri{\item Lo stesso errore ripetuto vale una volta sola.}
\end{boxcom}
\end{document}
